\documentclass[10pt,a4paper]{article}
\usepackage[spanish,es-nodecimaldot]{babel}
\usepackage[utf8]{inputenc}
\usepackage[T1]{fontenc}
\usepackage{lmodern}
\usepackage[top=1.1cm,bottom=1.1cm,left=1.5cm,right=1.5cm]{geometry}
\usepackage{amsmath,amssymb}

\setlength{\parindent}{0pt}
\setlength{\parskip}{2pt}
\pagestyle{empty}

\begin{document}

\textbf{Regresion multiple} $\longrightarrow$ \textit{A diferencia de la regresión simple, donde se analiza la relación de una única variable independiente con la dependiente ($Y$), la regresión múltiple incorpora dos o más variables independientes ($x_1, x_2, \ldots, x_k$) para explicar el comportamiento de $Y$. Esto se ajusta mucho mejor a los fenómenos prácticos, ya que los procesos complejos rara vez dependen de un solo factor. Desde una perspectiva geométrica, este cambio implica que dejamos de trabajar en un plano bidimensional buscando una recta de ajuste; en su lugar, operamos en un espacio tridimensional (o multidimensional), donde el objetivo matemático es estimar un plano (o hiperplano) de regresión.}

\end{document}
